\section{Estado actual y siguientes pasos}

\begin{frame}{Conclusiones hasta ahora}
  \begin{enumerate}
    \item El pipeline de datos y el grafo heterog\'eneo ya est\'an operativos y producen artefactos reutilizables.
    \item La variante V2 es, hasta ahora, la mejor configuraci\'on para combinar clasificaci\'on celular y predicci\'on de enlaces miRNA-gen.
    \item El proyecto ya genera resultados interpretables: saliencia por tipo celular, circuitos regulatorios y enriquecimiento funcional.
    \item Existen candidatos con inter\'es biol\'ogico directo para EM, pero necesitan validaci\'on externa y contraste con literatura/experimento.
  \end{enumerate}
\end{frame}

\begin{frame}{Limitaciones y riesgos actuales}
  \begin{itemize}
    \item La expresi\'on de miRNAs proviene de bulk, no de c\'elula \`unica; la resoluci\'on celular es indirecta.
    \item Parte del conocimiento miRNA-gen viene de predicci\'on computacional y no exclusivamente de validaci\'on experimental.
    \item La clasificaci\'on celular es muy fuerte; la pregunta dif\'icil y m\'as valiosa sigue siendo la causalidad regulatoria.
    \item Falta profundizar en diferencias MS vs control y estabilidad entre seeds o particiones.
  \end{itemize}
\end{frame}

\begin{frame}{Roadmap para la siguiente iteraci\'on de esta presentaci\'on}
  \small
  \begin{tabularx}{\textwidth}{>{\bfseries}p{3.2cm}X}
    \toprule
    Siguiente bloque & Acci\'on sugerida \\
    \midrule
    Datos & A\~nadir tablas de distribuci\'on MS vs control y recuentos por subtipo celular \\
    Resultados & Incorporar comparaci\'on por seed, por checkpoint y por ablation con narrativa homog\'enea \\
    Biolog\'ia & Resumir soporte bibliogr\'afico de miR-146a-3p, miR-140-5p y miR-23a-3p \\
    Validaci\'on & Integrar datasets externos o literatura curada para evaluar top circuits \\
    Visuales & Duplicar slides de red para Microglia, Monocyte y Oligodendrocyte \\
    \bottomrule
  \end{tabularx}
\end{frame}

\begin{frame}{Cierre}
  \begin{block}{Mensaje central}
    El proyecto ya pas\'o de integrar datos heterog\'eneos a producir hip\'otesis reguladoras concretas, cuantificables y visualizables por tipo celular en EM.
  \end{block}

  \vspace{0.8em}

  \centering
  \Large
  Siguiente meta: convertir los mejores circuitos en evidencia biol\'ogica defendible.
\end{frame}